\section{习题 2.4：平面、直线的距离与夹角（14 题）}\label{sec:24}
\tip{\textbf{本节四个核心公式：}\vspace{-2mm}
\begin{align*}
d(P,\pi)&=\frac{|Ax_0+By_0+Cz_0+D|}{\sqrt{A^2+B^2+C^2}},\\
d(P,l)&=\frac{|(P-P_0)\times u|}{|u|},\\
\cos\angle(\pi_1,\pi_2)&=\frac{|n_1\cdot n_2|}{|n_1||n_2|},\\
\sin\angle(l,\pi)&=\frac{|u\cdot n|}{|u||n|}.
\end{align*}
异面直线 $l_i=P_i+t_i u_i$ 的距离为
\[d(l_1,l_2)=\frac{|(P_2-P_1)\cdot(u_1\times u_2)|}{|u_1\times u_2|}.\]
}
\begin{center}
\begin{tikzpicture}[scale=.83,>=Stealth]
\fill[blue!9] (0,0)--(5.5,.4)--(7,2.4)--(1.5,2.0)--cycle;
\draw[blue!65] (0,0)--(5.5,.4)--(7,2.4)--(1.5,2.0)--cycle;
\coordinate (H) at (3.3,1.0);\coordinate (P) at (3.3,3.35);
\draw[->,very thick,red] (H)--(P) node[midway,right]{$d(P,\pi)$};
\fill (H) circle(1.4pt) node[below]{$H$};\fill (P) circle(1.4pt) node[above]{$P$};
\draw[thin] (H)--($(H)+(.35,.1)$)--($(H)+(.35,.45)$);
\node at (9.1,2) {$n=(A,B,C)$};
\node at (9.2,1.25) {$PH\parallel n$};
\end{tikzpicture}
\end{center}

\Q{1}{根据垂直或平行条件写出平面}
\begin{enumerate}[label=(\arabic*)]
\item 过 $P(-1,2,0)$，法向量 $n=(3,3,1)$，点法式直接得 $3x+3y+z-3=0$。
\item 过 $(-1,1,-10),(5,-2,11)$，垂直于 $2x+5y+z+1=0$。取两点差 $(6,-3,21)$，与给定法向量 $(2,5,1)$ 叉乘取法向量 $(3,-1,-1)$。\ans{$3x-y-z-6=0$。}
\item 平面垂直于 $x+y+z+4=0$，同时过两平面 $3x-4y-5z+1=0$ 与 $-2x+3y+2z-2=0$ 的交线。设平面族 $\lambda F_1+\mu F_2=0$，要求新法向量和 $(1,1,1)$ 点乘为零；可得\ans{$x-2y+z+3=0$。}
\item 过 $(2,0,-3)$ 并垂直于两平面 $x-2y+4z-7=0$、$3x+5y-2z+1=0$，法向量应平行两已知法向量的叉积 $(-16,14,11)$。\ans{$-16x+14y+11z+65=0$。}
\end{enumerate}
\ans{四小问与书后答案一致。}

\Q{2}{证明一个三阶行列式给出垂直于两平面的平面}
给出的行列式为
\[
\begin{vmatrix}x-x_0&y-y_0&z-z_0\\A_1&B_1&C_1\\A_2&B_2&C_2\end{vmatrix}=0.
\]
将第二、三行记作 $n_1,n_2$，该行列式等于 $(X-P_0)\cdot(n_1\times n_2)$。因此新平面的法向量 $n_1\times n_2$ 同时垂直于 $n_1,n_2$，于是新平面与原两平面互相垂直。
\warning{只有 $n_1,n_2$ 不平行时，叉积非零、上述行列式才真正定义一张平面。原题默认应排除两法向量成比例的退化情况。}

\Q{3}{根据正交条件写直线}
\begin{enumerate}[label=(\arabic*)]
\item 过 $P=(3,-2,1)$ 并垂直于平面 $3x+2y-3z+5=0$：方向就是 $(3,2,-3)$，\ans{$\displaystyle\frac{x-3}{3}=\frac{y+2}{2}=\frac{z-1}{-3}$。}
\item 过 $P=(1,1,-3)$，与过 $Q=(2,-1,2)$、方向 $u=(-1,2,4)$ 的直线正交并相交。垂足 $H=Q+tu$ 必须满足 $(H-P)\cdot u=0$；联立得到 $H$，联线 $PH$ 可写为\ans{$\begin{cases}-x+2y+4z+11=0,\\2x+y-3=0.\end{cases}$}
\item 过 $P=(0,1,-1)$，与交线 $3x+2y-5=0,2x-z+3=0$ 正交且相交。以该交线上的通点 $H$ 表示，使用 $\overrightarrow{PH}\cdot u=0$ 确定垂足，再将 $P,H$ 连成直线，\ans{$\begin{cases}18x+8y-3z-11=0,\\2x-3y+4z+7=0.\end{cases}$}
\end{enumerate}

\Q{4}{点到平面的距离}
(1) $P(2,1,0)$，平面 $3x-4y-5z+1=0$：
\[d=\frac{|6-4+1|}{\sqrt{9+16+25}}=\frac{3}{5\sqrt2}=\frac{3\sqrt2}{10}.\]
(2) $P(2,4,-1)$，平面 $x-z-1=0$：
\[d=\frac{|2-(-1)-1|}{\sqrt{1+1}}=\sqrt2.\]
\ans{$3\sqrt2/10$ 和 $\sqrt2$，均与原书一致。}

\Q{5}{两张平行平面的距离}
$2x-2y+z+5=0$ 与 $2x-2y+z-1=0$ 具有相同法向量 $(2,-2,1)$，故
\[d=\frac{|5-(-1)|}{\sqrt{2^2+(-2)^2+1^2}}=\frac63=2.\]
\ans{$d=2$。}

\Q{6}{四面体体积}
底面 $ABC$ 所在平面为 $3x-2y-6z-4=0$，面积为 $4$，顶点 $D(2,3,1)$。
\[h=d(D,ABC)=\frac{|6-6-6-4|}{\sqrt{9+4+36}}=\frac{10}{7}.\]
四面体体积 $V=\frac13Sh$，故
\ans{$\displaystyle V=\frac13\cdot4\cdot\frac{10}{7}=\frac{40}{21}$，与书后答案一致。}

\Q{7}{平面夹角、直线夹角和线面夹角}
\idea{平面夹角看法向量，直线夹角看方向向量，线面夹角用正弦而非余弦。若题目要求两面的\emph{锐}夹角，取绝对值。}
\begin{enumerate}[label=(\arabic*)]
\item $n_1=(3,-4,-5)$、$n_2=(4,1,-1)$，
\[n_1\cdot n_2=13,\quad |n_1||n_2|=30\Rightarrow\theta=\arccos\frac{13}{30}.\]
\item $u_1=(-1,1,2)$、$u_2=(-2,4,-3)$，内积 $2+4-6=0$，\ans{$\theta=\pi/2$。}
\item 两条直线为两平面交线，分别取两个法向量叉积为方向 $u_1,u_2$，整理 $|u_1\cdot u_2|/(|u_1||u_2|)$ 得\ans{$\displaystyle\theta=\arccos\frac{2\sqrt{22}}{11}$。}
\item $u=(2,1,-1)$、平面 $x-2y+4z-1=0$，$u\cdot n=2-2-4=-4$，\ans{$\displaystyle\theta=\arcsin\frac{2\sqrt{14}}{21}$。}
\item 从两张平面的交线求 $u=n_1\times n_2$，与 $2x-z-4=0$ 的法向量作正弦内积，\ans{$\displaystyle\theta=\arcsin\frac{\sqrt{70}}{14}$。}
\item 对 $\pi_1:3x-4y=0$、$\pi_2:y+z+1=0$、$\pi_3:x+2z+3=0$，先用 $u_{12}=n_1\times n_2$、$u_{13}=n_1\times n_3$。交线与第三平面的夹角及两交线的夹角分别为
\[
\arcsin\frac{2\sqrt{170}}{170},\qquad
\arccos\frac{31\sqrt{986}}{986}.
\]
\end{enumerate}
\ans{六小问与书后结果一致。}

\Q{8}{直线与平面平行时求参数及距离}
直线 $\frac{x+1}{t}=\frac{y-3}{1}=\frac{z+1}{2}$ 的方向向量 $u=(t,1,2)$，平面 $2x+4y-5z+1=0$ 的法向量 $n=(2,4,-5)$。
平行条件 $u\cdot n=2t+4-10=0$，所以 $t=3$。
取线上点 $P=(-1,3,-1)$，点面距离为
\[d=\frac{|2(-1)+4\cdot3-5(-1)+1|}{\sqrt{4+16+25}}=\frac{16}{3\sqrt5}=\frac{16\sqrt5}{15}.\]
\ans{$t=3$，$d=16\sqrt5/15$。}

\Q{9}{与给定平面距离恰好为 3 的点的轨迹}
给定平面 $4x-4y-2z-3=0$，法向量长度 $\sqrt{16+16+4}=6$。
\[\frac{|4x-4y-2z-3|}{6}=3 \iff 4x-4y-2z-3=\pm18.\]
\ans{两张平行平面 $4x-4y-2z-21=0$ 与 $4x-4y-2z+15=0$。}

\Q{10}{过 $z$ 轴并与给定平面夹角为 $60^\circ$}
任何过 $z$ 轴的平面都可写为 $Ax+By=0$（$(A,B)\ne0$），其法向量 $n_1=(A,B,0)$。已知 $\pi:2x+y-\sqrt5\,z-1=0$ 的法向量 $n_2=(2,1,-\sqrt5)$。
由
\[\frac{|2A+B|}{\sqrt{10}\sqrt{A^2+B^2}}=\cos60^\circ=\frac12\]
整理出关于 $A:B$ 的两个方向比例，最终两张平面为
\ans{$x+3y=0$ 和 $3x-y=0$，与书后答案一致。}

\Q{11}{到两非平行平面距离相等的轨迹}
两平面为 $3x-2y-6z-4=0$ 与 $2x+2y-z+5=0$，法向量长度分别为 $7$ 和 $3$。故
\[\frac{|3x-2y-6z-4|}{7}=\frac{|2x+2y-z+5|}{3}.\]
去掉绝对值分正负两支：$3F_1-7F_2=0$ 和 $3F_1+7F_2=0$。
\ans{$23x+8y-25z+23=0$ 及 $5x+20y+11z+47=0$。它们是这两张平面的两个角平分面。}

\Q{12}{点到直线的距离}
\idea{把直线交线化为 $P_0+tu$，再用 $|(P-P_0)\times u|/|u|$。}
\begin{enumerate}[label=(\arabic*)]
\item 点 $P=(1,0,2)$，直线为 $3x-2y-1=0,\ y-z-2=0$。取线上点 $P_0=(1,1,-1)$（可复验两式），方向 $u=n_1\times n_2$，于是
\[d=\frac{|(P-P_0)\times u|}{|u|}=\frac{\sqrt{1012}}{11}.\]
\item 点 $P=(3,10,-1)$，直线为 $\frac{x-1}{2}=\frac{y-5}{10}=\frac{z+1}{-3}$。取 $P_0=(1,5,-1)$，$u=(2,10,-3)$，直接代入得
\[d=\frac{|(2,5,0)\times(2,10,-3)|}{\sqrt{113}}=\frac{19\sqrt{113}}{113}.\]
\end{enumerate}
\ans{结果与原书一致。}

\Q{13}{两条直线之间的距离}
先检查 $u_1\times u_2$：若非零且两线异面，用混合积除以方向叉积模；若两线平行，再取基点差对共同方向求垂距。逐小题求出基点和方向，按实际情况代入得
\ans{$\displaystyle (1)\ 2\sqrt3,\qquad (2)\ \frac{3\sqrt{41}}{41},\qquad(3)\ \frac{\sqrt6}{6}$。}
\tip{求两线距离先判	extbf{平行或异面}；两种公式不同。平行线换基点不会改变叉积垂距，异面线换基点也不会改变混合积垂距。}

\Q{14}{异面直线的距离与公垂线}
\idea{异面两线的公垂线方向 $w=u_1\times u_2$，先用 $d=|\Delta P\cdot w|/|w|$ 算距离；再在两线上取动点 $P_1+s u_1,P_2+t u_2$，令二者连线同时垂直 $u_1,u_2$，解出 $s,t$，就得到公垂线。}
\begin{enumerate}[label=(\arabic*)]
\item 对
\[\ell_1:\begin{cases}3x+y-3=0,\\y+z=0,\end{cases}\qquad
\ell_2:\begin{cases}x+z=0,\\x-2y=0,\end{cases}\]
先解 $P_i,u_i$，分别取动点 $X_i$，由 $(X_2-X_1)\cdot u_1=(X_2-X_1)\cdot u_2=0$ 求两垂足。代入得
\ans{$d=3\sqrt{122}/122$，公垂线 $\begin{cases}45x-2y-17z-45=0,\\23x-20y+13z=0.\end{cases}$}
\item 对题设两条点向式直线 $\ell_1:(x-1)/(-1)=y/1=z/0$、$\ell_2:(x-1)/2=y/(-1)=(z-2)/2$，同样计算得到
\ans{$d=2/3$，公垂线 $\begin{cases}x+y+4z-1=0,\\x-2y-2z+3=0.\end{cases}$}
\end{enumerate}
